% !TeX program = pdflatex
% =============================================================================
% Aufgabenblatt 2: Kraftaddition & Kraftzerlegung
% UZH Praktikum I -- Sara Romer Urban, Kantonsschule im Lee, HS2026
% Klasse: 2e (02.09.2026)
%
% Reine Uebungsaufgaben zu Kraftaddition/-zerlegung (Sara Romers Vorgabe,
% siehe kraefte-addieren-zerlegen-lektionsplan.tex Abschnitt 5), alle
% grafisch, kein einleitender Erklaerkasten. Bei allen bildbasierten
% Aufgaben (Schiff, Seilbahn, Seilzug) zeichnen die SuS direkt in die
% jeweilige Abbildung hinein -- kein separates Feinraster zum Uebertragen
% mehr. Fuenf Aufgaben: Schiff (Addition zweier gegebener, konkreter
% Schlepperkraefte 40 kN/40 kN bei selbst gewaehltem Kraftmassstab, danach
% aus der grafisch bestimmten Resultierenden und der Schiffsmasse die
% Beschleunigung berechnen), Seilbahn (Zerlegung der Gewichtskraft im
% Gleichgewicht, rein graphisch ohne Kraftmassstab noetig), Seilzug mit drei
% Gewichten (Addition zweier gegebener, konkreter Seilkraefte 180 N/380 N
% bei selbst gewaehltem Kraftmassstab, um auf die Gewichtskraft des
% mittleren, unbekannten Gewichts zu kommen -- das fruehere
% Seilziehen-Beispiel steht jetzt stattdessen in der Vektor-Theorie von
% arbeitsblatt3-kraefte-addieren-zerlegen.tex), ein reines
% Gitter-Konstruktions-Beispiel (Aufgabe 4: drei vorgegebene Kraftpfeile,
% Massstab gegeben, Kopf-an-Schwanz-Addition, Betrag der Resultierenden
% ablesen), sowie neu Aufgabe 5 (Lampe an durchhaengendem Seil: zwei exakt
% per Trigonometrie konstruierte, bewusst gross skalierte TikZ-Skizzen mit
% Durchhangwinkel 40°/20°, je rein grafische Kraftzerlegung mit selbst
% gewaehltem Kraftmassstab -- Loesungen pruefen nur per
% Kopf-an-Schwanz-Gleichgewicht, keine Sinus-Formel --, Vergleich der beiden
% Seilspannkraefte, und als Grenzfall-Ueberlegung die rein geometrische
% Frage, wie lang die Seilspannkraefte bei einem Winkel von ca. 0.1°
% gezeichnet werden muessten: zeigt ueber die Parallelogramm-Konstruktion
% (ohne Formel), warum ein Seil nie exakt waagrecht gespannt werden kann).
% Die fruehere Hundeschlitten-Aufgabe wurde gestrichen. Theorie, Lernziele
% und alles im Plenum Behandelte stehen in jenem Arbeitsblatt.
%
% Quellen: lernziele.md, kraefte-addieren-zerlegen-lektionsplan.tex (dieser
% Ordner), Bilder/schiff_aufgabe(.png/_loesung.png),
% Bilder/seilbahn_aufgabe(.png/_loesung.png),
% Bilder/kraftzerlegung_seilzug_aufgabe.png.
% =============================================================================
\documentclass[addpoints,11pt,a4paper]{exam}

\usepackage[T1]{fontenc}
\usepackage[utf8]{inputenc}
\usepackage[ngerman]{babel}
\usepackage{lmodern}
\usepackage[a4paper,margin=2.2cm,headheight=14pt]{geometry}
\usepackage{amsmath}
\usepackage{siunitx}
% Hausstil: zusammengesetzte Einheiten immer als echter Bruch (\frac), nicht
% als Inline-Schraegstrich oder \tfrac -- gilt global fuer jedes \si/\SI.
\sisetup{per-mode=fraction,fraction-command=\frac}
\usepackage{array,booktabs,tabularx,multirow}
\usepackage{enumitem}
\usepackage{xcolor}
\usepackage{colortbl}
\usepackage{tikz}
\usepackage{physics}
\usepackage[most]{tcolorbox}
\usepackage{lastpage}
\usepackage{graphicx}
\usepackage{hyperref}

\usetikzlibrary{arrows.meta,calc,angles,quotes,decorations.pathmorphing}

\usepackage{setspace}
\usepackage{needspace}
\setlength{\parindent}{0pt}
\setlength{\parskip}{0.6em}
\setstretch{1.08}
\setlist[itemize]{itemsep=0.4em}

% -----------------------------
% Farben (Corporate-Farbschema Physik KFR)
% -----------------------------
\definecolor{titleblue}{HTML}{183A6B}
\definecolor{accentblue}{HTML}{2E6F9E}
\definecolor{lightblue}{HTML}{EAF4FB}
\definecolor{lightgray}{HTML}{F4F6F8}
\definecolor{darkgray}{HTML}{4A4A4A}
\definecolor{accentorange}{HTML}{D9720C}
\definecolor{accentpurple}{HTML}{7B2D8E}
% Hinweisbox: Orange (accentorange), fuer wichtige Klarstellungen.
\definecolor{lightorange}{HTML}{FCEEE0}
% Formelbox: Violett (accentpurple), fuer die zentralen Merksaetze.
\definecolor{lightpurple}{HTML}{F3EAF7}

% Links (hyperref): farblich ans Corporate-Farbschema angeglichen.
\hypersetup{
  colorlinks=true,
  urlcolor=accentblue,
  linkcolor=accentblue,
  pdftitle={Aufgabenblatt 2: Kraftaddition \& Kraftzerlegung},
  pdfauthor={Loris Keller}
}

% -----------------------------
% Spaltentypen
% -----------------------------
\newcolumntype{Y}{>{\centering\arraybackslash}X}
\newcolumntype{P}[1]{>{\raggedright\arraybackslash}p{#1}}

% -----------------------------
% Boxen
% -----------------------------
\tcbset{before skip=10pt, after skip=10pt}
\newtcolorbox{titlebox}{
  enhanced,
  colback=titleblue,
  colframe=titleblue,
  boxrule=0pt,
  arc=3mm,
  left=3mm,right=3mm,top=3mm,bottom=3mm
}

\newlength{\aufgabenindent}
\settowidth{\aufgabenindent}{10.\hskip\labelsep}

% Praktikum I: Lernziele/Einleitung/Theorie/Aufgaben werden NICHT mehr
% farbig gerahmt (nur Formel-, Hinweis- und Antwortbox bleiben Boxen) --
% siehe institutionen/UZH-Praktikum-I/CLAUDE.md, Abschnitt "Boxen-Layout".
\newenvironment{infobox}{%
  \par\addvspace{10pt}\noindent
}{%
  \par\addvspace{10pt}
}

\newenvironment{theoriebox}[1][Theorie]{%
  \par\addvspace{10pt}\noindent\textbf{#1}\par\smallskip\noindent
}{%
  \par\addvspace{10pt}
}

% taskbox/groupbox stehen in questions VOR dem ersten \question (Bilder,
% Fliesstext) -- exam.cls' questions-Liste braucht dafuer eine echte,
% eigenstaendige Box (sonst "Something's wrong--perhaps a missing \item",
% weil z.B. ein \begin{center} vor dem ersten \item der Aussenliste steht).
% Deshalb bleiben es tcolorboxen, nur unsichtbar gemacht (keine Farbe/Rahmen,
% kein Padding) statt sie durch reine Absaetze zu ersetzen.
\newtcolorbox{taskbox}[1][]{
  enhanced, breakable,
  colback=white, colframe=white, boxrule=0pt,
  left=0mm,right=0mm,top=0mm,bottom=0mm,
  before skip=10pt, after skip=10pt,
  before upper={%
    \def\tmpboxtitle{#1}%
    \ifx\tmpboxtitle\empty\else\noindent\textbf{#1}\par\smallskip\fi
  }
}

\newtcolorbox{groupbox}[1][Gruppenarbeit]{
  enhanced, breakable,
  colback=white, colframe=white, boxrule=0pt,
  left=0mm,right=0mm,top=0mm,bottom=0mm,
  before skip=10pt, after skip=10pt,
  before upper={\noindent\textbf{#1}\par\smallskip}
}

% beispielbox: fuer die Einleitung -- wie die anderen Inhaltsboxen ungerahmt.
\newenvironment{beispielbox}[1][Beispiel]{%
  \par\addvspace{10pt}\noindent\textbf{#1}\par\smallskip\noindent
}{%
  \par\addvspace{10pt}
}

% hinweisbox: Orange, fuer wichtige Klarstellungen/Verwechslungswarnungen.
\newtcolorbox{hinweisbox}[1][Hinweis]{
  enhanced,
  breakable,
  colback=lightorange,
  colframe=accentorange,
  title={#1},
  fonttitle=\bfseries,
  boxrule=0.7pt,
  arc=2mm,
  left=5mm,right=5mm,top=2mm,bottom=2mm
}

% formelbox: Violett, um die zentralen Merksaetze dieser Einheit hervorzuheben.
\newtcolorbox{formelbox}[1][Merksatz]{
  enhanced,
  breakable,
  colback=lightpurple,
  colframe=accentpurple,
  title={#1},
  fonttitle=\bfseries,
  boxrule=0.9pt,
  arc=2mm,
  left=5mm,right=5mm,top=2mm,bottom=2mm
}

\newtcolorbox{answerbox}{
  breakable,
  colback=white,
  colframe=gray!55,
  boxrule=0.5pt,
  arc=1.5mm,
  left=2mm,right=2mm,top=2mm,bottom=2mm
}

% Marke: kleines farbiges Inline-Label fuer Teilschritte/Ergebnis-Callouts.
\newtcbox{\Marke}[1][accentpurple]{
  nobeforeafter,
  colback=#1,
  colframe=#1,
  coltext=white,
  boxrule=0pt,
  arc=2pt,
  left=5pt,right=5pt,top=2pt,bottom=2pt,
  fontupper=\bfseries\sffamily\small
}

% --- Feinraster-Helfer fuer alle grafischen Konstruktionsaufgaben ----------
% Feld mit feinem Punkt-/Liniengitter (0.5 Kaestchen Nebenlinien, 1 Kaestchen
% Hauptlinien), sodass bei Massstab 1 Kaestchen = 1 N jede Kraft exakt und
% nachzaehlbar eingezeichnet werden kann. #1#2 = linke untere Ecke, #3#4 =
% rechte obere Ecke.
\newcommand{\Feinraster}[4]{%
  \draw[gray!22,step=0.5] (#1,#2) grid (#3,#4);
  \draw[gray!55,step=1] (#1,#2) grid (#3,#4);
}

% Bild-Helfer: zentriertes Bild mit spuerbarem Abstand zum umgebenden Text
% (statt nur des normalen Absatzabstands). #1 = Breite, #2 = Dateipfad.
\newcommand{\Bild}[2]{%
  \par\vspace{0.6cm}
  \begin{center}
    \includegraphics[width=#1]{#2}
  \end{center}
  \vspace{0.6cm}
}

% --- Loesungen ein-/ausblenden -----------------------------------------------
\noprintanswers
%\printanswers

\pointformat{}
\renewcommand{\solutiontitle}{\noindent\color{titleblue}\textbf{L\"osung:}\enspace}

\pagestyle{headandfoot}
\cfoot{\textcolor{darkgray}{Seite \thepage\ von \pageref{LastPage}}}

\newcommand{\Kopfzeile}[4]{%
  \begin{titlebox}
    {\color{white}\sffamily\bfseries\Large #4}
  \end{titlebox}
  \vspace{0.5em}
  \noindent\textbf{Klasse:}~#1 \hfill \textbf{Datum:}~#2 \hfill \textbf{Thema:}~#3
    \hfill \textbf{Lehrperson:}~Loris Keller
  \vspace{0.8em}
  \lhead{\textcolor{darkgray}{#3}}
  \rhead{\textcolor{darkgray}{#4}}
}

\newlength{\antwortraumlen}
\newcommand{\Antwortraum}[1]{%
  \pgfmathsetlength{\antwortraumlen}{#1*2.5cm}%
  \begin{answerbox}
    \rule{0pt}{\antwortraumlen}%
  \end{answerbox}%
}

% Werte fuer Kopfzeile zentral setzen
\newcommand{\Klasse}{2e}
\newcommand{\Datum}{02.09.2026}
\newcommand{\Thema}{Kraftaddition \& Kraftzerlegung}
\newcommand{\ArbeitsblattTitel}{Aufgabenblatt 2: Kraftaddition \& Kraftzerlegung}

\begin{document}

\Kopfzeile{\Klasse}{\Datum}{\Thema}{\ArbeitsblattTitel}

% =============================================================================
% Aufgabe 1: Schiff (Addition zweier Schlepperkraefte)
% =============================================================================
\begin{questions}
\begin{taskbox}[Aufgabe 1: Zwei Schlepper ziehen ein Schiff]
\question[7] Ein Containerschiff wird gleichzeitig von zwei Schleppern mit
den Kräften $\vec F_1=\SI{40}{\kilo\newton}$ und $\vec F_2=\SI{40}{\kilo\newton}$
gezogen:

\Bild{15.5cm}{Bilder/schiff_aufgabe.png}

\begin{parts}
\part[4] Zeichnen Sie direkt in die Skizze oben: Wählen Sie einen selbst
gewählten Kraftmassstab, und zeichnen Sie $\vec F_1$ und $\vec F_2$ mit den
entsprechenden Längen entlang der eingezeichneten Schlepprichtungen ein
(Winkel mit dem Geodreieck von der Abbildung abmessen). Bestimmen Sie
anschliessend grafisch (Kräfteparallelogramm) den Betrag der
Resultierenden $\vec F_R=\vec F_1+\vec F_2$, mit der sich das Schiff
tatsächlich bewegt.

\begin{solution}
\Bild{14cm}{Bilder/schiff_aufgabe_loesung.png}
Mit gleich grossen Kräften $F_1=F_2=\SI{40}{\kilo\newton}$ liegt die
Resultierende genau auf der Schiffsachse; ihr Betrag ergibt sich (je nach
gemessenem Winkel) zu ungefähr $F_R\approx\SI{69}{\kilo\newton}$.
\end{solution}

\part[3] Das Containerschiff hat eine Masse von $m=\SI{2000}{\tonne}$.
Nutzen Sie Ihr Ergebnis aus (a), um die Beschleunigung $a$ zu berechnen,
mit der das Schiff aus dem Stillstand heraus beschleunigt wird.

\vspace{7.5cm}

\begin{solution}
Newtons zweites Gesetz nach $a$ auflösen:
\[
F_R = m\cdot a \quad\Longrightarrow\quad a = \frac{F_R}{m}
\]
Werte einsetzen (mit $F_R\approx\SI{69}{\kilo\newton}=\SI{69000}{\newton}$
und $m=\SI{2000}{\tonne}=\SI{2000000}{\kilogram}$):
\[
a = \frac{\SI{69000}{\newton}}{\SI{2000000}{\kilogram}}
  \approx \SI{0.035}{\meter\per\second\squared}
\]
\end{solution}
\end{parts}
\end{taskbox}
\end{questions}

% =============================================================================
% Aufgabe 2: Seilbahn (Zerlegung, Gleichgewicht)
% =============================================================================
\newpage
\begin{questions}
\begin{taskbox}[Aufgabe 2: Seilbahngondel im Gleichgewicht]
\question[5] Eine Seilbahngondel hängt in Ruhe am Tragseil. Die Gewichtskraft
$\vec F_g$ der Gondel samt Passagieren muss deshalb vollständig von den
beiden Seilabschnitten aufgenommen werden: Es gilt
$\vec F_{S1}+\vec F_{S2}+\vec F_g=\vec F_{\text{res}}=\vec 0$.

Zeichnen Sie direkt in die Skizze unten: Messen Sie die eingezeichneten
Seilrichtungen mit dem Geodreieck ab, und zerlegen Sie $\vec F_g$ grafisch
(ohne Kraftmassstab) in die beiden Seilkraftkomponenten $\vec F_{S1}$ und
$\vec F_{S2}$.

\vspace{2cm}
\Bild{15.5cm}{Bilder/seilbahn_aufgabe.png}

\begin{solution}
\Bild{15cm}{Bilder/seilbahn_aufgabe_loesung.png}
\end{solution}
\end{taskbox}
\end{questions}

% =============================================================================
% Aufgabe 3: Kraefte am Seilzug (Addition zweier gegebener Seilkraefte)
% =============================================================================
\newpage
\begin{questions}
\begin{taskbox}[Aufgabe 3: Kräfte am Seilzug]
\question[5] An einem Seilzug mit zwei Rollen hängen zwei Gewichte mit den
Gewichtskräften $F_1=\SI{180}{\newton}$ und $F_2=\SI{380}{\newton}$. Am
Knotenpunkt in der Mitte hängt ein drittes, noch unbekanntes Gewicht
senkrecht nach unten:

Zeichnen Sie direkt in die Skizze unten: Wählen Sie einen selbst gewählten
Kraftmassstab, und zeichnen Sie $\vec F_1$ und $\vec F_2$ mit den
entsprechenden Längen entlang der eingezeichneten Seilrichtungen ein
(Winkel mit dem Geodreieck von der Abbildung abmessen). Bestimmen Sie
anschliessend ebenfalls direkt in der Skizze grafisch
(Kräfteparallelogramm) die Resultierende $\vec F_1+\vec F_2$: Ihr Betrag
entspricht der Gewichtskraft des mittleren, unbekannten Gewichts.

\Bild{15cm}{Bilder/kraftzerlegung_seilzug_aufgabe.png}

\begin{solution}
Damit der Knoten in Ruhe bleibt, muss die Resultierende $\vec F_1+\vec F_2$
senkrecht nach unten zeigen. Das ist eine gute Kontrolle für die eigene
Konstruktion. Ihr Betrag ergibt (je nach gemessenem Winkel) die gesuchte
Gewichtskraft des mittleren Gewichts, ungefähr $\SI{390}{\newton}$.
\end{solution}
\end{taskbox}
\end{questions}

% =============================================================================
% Aufgabe 4: Mehrere Kraefte im Gitter addieren (Gitter + Massstab gegeben)
% =============================================================================
\newpage
\begin{questions}
\begin{taskbox}[Aufgabe 4: Mehrere Kräfte im Gitter addieren]
\question[5] Am Punkt $A$ greifen drei Kräfte $\vec F_1$, $\vec F_2$,
$\vec F_3$ an (Massstab: $1$ Kästchen $\hat=$ $\SI{4}{\newton}$):

\begin{center}
\begin{tikzpicture}[>=Latex,scale=1]
  \Feinraster{-3}{-4}{9}{7}
  \coordinate (A) at (0,0);
  \draw[->,thick,titleblue] (A) -- (3,0) node[midway,above]{\scriptsize $\vec F_1$};
  \draw[->,thick,accentorange] (A) -- (0,4) node[midway,left]{\scriptsize $\vec F_2$};
  \draw[->,thick,accentpurple] (A) -- (1,-1) node[below right=1pt]{\scriptsize $\vec F_3$};
  \node[below left] at (A) {$A$};
\end{tikzpicture}
\end{center}

Konstruieren Sie mit dem Kopf-an-Schwanz-Verfahren direkt im Gitter die
Resultierende $\vec F_{\text{ges}}=\vec F_1+\vec F_2+\vec F_3$. Lesen Sie
deren Länge in Kästchen ab, und bestimmen Sie mit dem angegebenen Massstab
den Betrag $F_{\text{ges}}$ in Newton.

\begin{solution}
\begin{center}
\begin{tikzpicture}[>=Latex,scale=1]
  \Feinraster{-3}{-4}{9}{7}
  \coordinate (A) at (0,0);
  \coordinate (P1) at (3,0);
  \coordinate (P2) at (3,4);
  \coordinate (P3) at (4,3);
  \draw[->,thick,titleblue] (A) -- (P1) node[midway,above]{\scriptsize $\vec F_1$};
  \draw[->,thick,accentorange] (P1) -- (P2) node[midway,right]{\scriptsize $\vec F_2$};
  \draw[->,thick,accentpurple] (P2) -- (P3) node[midway,above]{\scriptsize $\vec F_3$};
  \draw[->,ultra thick,darkgray] (A) -- (P3) node[midway,below right]{\scriptsize $\vec F_{\text{ges}}$, $5$ Kästchen};
\end{tikzpicture}
\end{center}
Die Resultierende ist $5$ Kästchen lang, also
$F_{\text{ges}}=5\cdot\SI{4}{\newton}=\SI{20}{\newton}$.
\end{solution}
\end{taskbox}
\end{questions}

% =============================================================================
% Aufgabe 5: Lampe an durchhaengendem Seil (Grenzfall kleiner Winkel)
% =============================================================================
\newpage
\begin{questions}
\begin{taskbox}[Aufgabe 5: Eine Lampe am gespannten Seil]
\question[7] Eine Lampe mit der Gewichtskraft $F_g=\SI{20}{\newton}$ hängt in
der Mitte eines Seils, das zwischen zwei gegenüberliegenden Wänden gespannt
ist. Durch das Gewicht der Lampe hängt das Seil auf beiden Seiten leicht
durch und bildet dort mit der Horizontalen jeweils den Winkel $\alpha$:

\begin{parts}
\part[2] \textbf{Skizze 1:} Hier beträgt der Winkel $\alpha_1=\SI{40}{\degree}$.

\vspace{1.5cm}
\begin{center}
\begin{tikzpicture}[>=Latex,scale=1.4]
  \def\halfwidth{4}
  \def\winkelgrad{40}
  \pgfmathsetmacro{\h}{\halfwidth*tan(\winkelgrad)}
  \coordinate (L) at (-\halfwidth,0);
  \coordinate (R) at (\halfwidth,0);
  \coordinate (P) at (0,-\h);
  \draw[thick,darkgray] (L) -- (P) -- (R);
  \fill[darkgray] (L) circle (2.5pt);
  \fill[darkgray] (R) circle (2.5pt);
  \fill[black] (P) circle (2.5pt);
  \draw[->,thick,accentorange] (P) -- ++(0,-1.2) node[right]{$\vec F_g$};
  \pic[draw=accentblue,angle radius=0.9cm] {angle=P--L--R};
  \pic[draw=accentblue,angle radius=0.9cm] {angle=L--R--P};
  \pgfmathsetmacro{\halfang}{\winkelgrad/2}
  \node[accentblue] at ($(L)+({1.3*cos(\halfang)},{-1.3*sin(\halfang)})$) {$\alpha_1$};
  \node[accentblue] at ($(R)+({-1.3*cos(\halfang)},{-1.3*sin(\halfang)})$) {$\alpha_1$};
  \node[left] at (L) {Wand};
  \node[right] at (R) {Wand};
\end{tikzpicture}
\end{center}

Zeichnen Sie direkt in die Skizze: Wählen Sie einen selbst gewählten
Kraftmassstab, zeichnen Sie $\vec F_g$ nach unten ein, und zerlegen Sie sie
grafisch (Kräfteparallelogramm) in die beiden Seilspannkräfte entlang der
eingezeichneten Seilrichtungen. Bestimmen Sie so den Betrag der
Seilspannkraft $F_{S,1}$.

\begin{solution}
Als Kontrolle für die grafische Konstruktion: Weil die Skizze symmetrisch
ist, sind die beiden Seilspannkräfte gleich lang. Addiert man sie zusammen
mit $\vec F_g$ Kopf an Schwanz, muss sich exakt $\vec 0$ ergeben, denn die
Lampe hängt in Ruhe (Kräftegleichgewicht).
\end{solution}

\part[2] \textbf{Skizze 2:} Dieselbe Lampe hängt nun straffer, sodass der
Winkel nur noch halb so gross ist wie in Skizze 1: $\alpha_2=\SI{20}{\degree}$.

\vspace{1.5cm}
\begin{center}
\begin{tikzpicture}[>=Latex,scale=1.4]
  \def\halfwidth{4}
  \def\winkelgrad{20}
  \pgfmathsetmacro{\h}{\halfwidth*tan(\winkelgrad)}
  \coordinate (L) at (-\halfwidth,0);
  \coordinate (R) at (\halfwidth,0);
  \coordinate (P) at (0,-\h);
  \draw[thick,darkgray] (L) -- (P) -- (R);
  \fill[darkgray] (L) circle (2.5pt);
  \fill[darkgray] (R) circle (2.5pt);
  \fill[black] (P) circle (2.5pt);
  \draw[->,thick,accentorange] (P) -- ++(0,-1.2) node[right]{$\vec F_g$};
  \pic[draw=accentblue,angle radius=0.9cm] {angle=P--L--R};
  \pic[draw=accentblue,angle radius=0.9cm] {angle=L--R--P};
  \pgfmathsetmacro{\halfang}{\winkelgrad/2}
  \node[accentblue] at ($(L)+({1.3*cos(\halfang)},{-1.3*sin(\halfang)})$) {$\alpha_2$};
  \node[accentblue] at ($(R)+({-1.3*cos(\halfang)},{-1.3*sin(\halfang)})$) {$\alpha_2$};
  \node[left] at (L) {Wand};
  \node[right] at (R) {Wand};
\end{tikzpicture}
\end{center}

Wiederholen Sie die Konstruktion aus (a) für diese Skizze, und bestimmen Sie
so den Betrag der Seilspannkraft $F_{S,2}$.

\begin{solution}
Als Kontrolle wie in (a): Auch hier sind die beiden Seilspannkräfte gleich
lang, und zusammen mit $\vec F_g$ Kopf an Schwanz addiert muss sich wieder
exakt $\vec 0$ ergeben.
\end{solution}

\part[1] Vergleichen Sie $F_{S,1}$ und $F_{S,2}$: Der Winkel in Skizze 2 ist
nur halb so gross wie in Skizze 1. Was fällt Ihnen beim Vergleich der
beiden Seilspannkräfte auf?

\Antwortraum{1}

\begin{solution}
Obwohl der Winkel halbiert wurde, hat sich die gemessene Seilspannkraft
nicht auch einfach verdoppelt: $F_{S,2}$ ist deutlich, aber nicht exakt
doppelt so gross wie $F_{S,1}$. Der Zusammenhang zwischen Winkel und
Seilspannkraft ist also nicht proportional. Je kleiner der Winkel, desto
stärker wächst die Seilspannkraft an (siehe Teilaufgabe (d)).
\end{solution}

\part[2] Stellen Sie sich vor, Sie versuchen, die Konstruktion aus (a) und
(b) mit einem winzigen Winkel von $\alpha\approx\SI{0.1}{\degree}$ zu
wiederholen (das Seil hängt dann fast gar nicht mehr durch, es ist fast
perfekt waagrecht). Überlegen Sie anhand des Kräfteparallelogramms: Wie
lang müssten die beiden Seilspannkräfte in dieser Konstruktion werden?
Erklären Sie damit, warum ein Seil nie exakt waagrecht gespannt werden
kann.

\Antwortraum{2}

\begin{solution}
Bei einem derart kleinen Winkel verlaufen beide Seilrichtungen praktisch
waagrecht, also fast genau parallel zueinander. Um beim
Kräfteparallelogramm von der Spitze bzw. vom Fusspunkt des kurzen Pfeils
$\vec F_g$ aus Parallelen zu diesen beiden fast gleichen Richtungen zu
ziehen, bis sie sich treffen, müsste man die beiden Seilspannkräfte extrem
lang zeichnen, viel länger, als auf ein Blatt Papier passt. Je kleiner der
Winkel wird, desto länger müssten diese Pfeile werden. Im Grenzfall eines
exakt waagrechten Seils ($\alpha=\SI{0}{\degree}$) wären die beiden
Seilrichtungen exakt parallel: Die Konstruktionslinien würden sich nie mehr
treffen, die Seilspannkraft müsste unendlich gross sein. Ein reales Seil
kann aber nur eine endliche Kraft aufbringen, bevor es reisst. Deshalb
hängt jedes gespannte Seil, an dem ein Gewicht hängt, immer ein kleines
bisschen durch; ein perfekt waagrecht gespanntes Seil ist physikalisch
unmöglich.
\end{solution}
\end{parts}
\end{taskbox}
\end{questions}

\end{document}
